(1) 自由下落满足 \(y(t)=h-gt^2/2\) 直到 \(t_0=\sqrt{2h/g}\)，碰前速度 \(-gt_0=-\sqrt{2gh}\)。理想墙使速度瞬时反号，回升到 h 再静止所需同样时间。因此完整往复周期 \(T=2t_0=2\sqrt{2h/g}\)。

(2) 周期平均值可只在下降半段积分并乘以2：
\[
\langle U\rangle=\frac1{t_0}\int_0^{t_0}mg(h-gt^2/2)\,dt
=mg(h-gt_0^2/6)=\frac23mgh.
\]
总能量恒为 \(mgh\)，所以 \(\langle K\rangle=mgh/3\)。直接积分 \(K=mg^2t^2/2\) 也得到同一数值。

(3) 对粒子在光滑运动区的 \(G=yp\)，有 \(\dot G=2K-yU'(y)=2K-mgy\)。墙碰撞处 y=0，故冲量对 G 的积分跳跃为零；一周期 G 返回原值，所以 \(2\langle K\rangle=\langle yU'(y)\rangle=\langle U\rangle\)。与平均值 \(2(mgh/3)=2mgh/3\) 相符。

(4) 对下降阶段 \(p=-mgt\)；从碰后开始令局部时间 \(s=t-t_0\)，则 \(p=m\sqrt{2gh}-mgs\)，其中 \(0<s<t_0\)。碰撞瞬间 \(p^+=-p^-\)，动能 \(p^2/(2m)\) 不变，势能 \(mgy=0\)，总能量守恒。
